\documentclass[11pt]{article}
% Shared preamble for the Carry-Adder Wall paper series.
% Each paper does:  % Shared preamble for the Carry-Adder Wall paper series.
% Each paper does:  % Shared preamble for the Carry-Adder Wall paper series.
% Each paper does:  \input{../../shared/preamble}
\usepackage[margin=1in]{geometry}
\usepackage{amsmath, amssymb, amsthm}
\usepackage{microtype}
\usepackage{graphicx}
\usepackage[hidelinks]{hyperref}

\newtheorem{theorem}{Theorem}
\newtheorem{corollary}{Corollary}
\newtheorem{lemma}{Lemma}
\theoremstyle{remark}
\newtheorem*{remark}{Remark}

\newcommand{\Mhash}{\mathsf{M}}
\newcommand{\SHA}{\mathrm{SHA\text{-}256}}
\newcommand{\SHAd}{\mathrm{SHA\text{-}256d}}
\newcommand{\MAJ}{\mathrm{MAJ}}
\newcommand{\Ch}{\mathrm{Ch}}
\newcommand{\Maj}{\mathrm{Maj}}
\newcommand{\bigO}{\mathcal{O}}

\usepackage[margin=1in]{geometry}
\usepackage{amsmath, amssymb, amsthm}
\usepackage{microtype}
\usepackage{graphicx}
\usepackage[hidelinks]{hyperref}

\newtheorem{theorem}{Theorem}
\newtheorem{corollary}{Corollary}
\newtheorem{lemma}{Lemma}
\theoremstyle{remark}
\newtheorem*{remark}{Remark}

\newcommand{\Mhash}{\mathsf{M}}
\newcommand{\SHA}{\mathrm{SHA\text{-}256}}
\newcommand{\SHAd}{\mathrm{SHA\text{-}256d}}
\newcommand{\MAJ}{\mathrm{MAJ}}
\newcommand{\Ch}{\mathrm{Ch}}
\newcommand{\Maj}{\mathrm{Maj}}
\newcommand{\bigO}{\mathcal{O}}

\usepackage[margin=1in]{geometry}
\usepackage{amsmath, amssymb, amsthm}
\usepackage{microtype}
\usepackage{graphicx}
\usepackage[hidelinks]{hyperref}

\newtheorem{theorem}{Theorem}
\newtheorem{corollary}{Corollary}
\newtheorem{lemma}{Lemma}
\theoremstyle{remark}
\newtheorem*{remark}{Remark}

\newcommand{\Mhash}{\mathsf{M}}
\newcommand{\SHA}{\mathrm{SHA\text{-}256}}
\newcommand{\SHAd}{\mathrm{SHA\text{-}256d}}
\newcommand{\MAJ}{\mathrm{MAJ}}
\newcommand{\Ch}{\mathrm{Ch}}
\newcommand{\Maj}{\mathrm{Maj}}
\newcommand{\bigO}{\mathcal{O}}


\title{\textbf{The Amortization Envelope:\\
Where a SHA-256d Mining Advantage Can and Cannot Live}}
\author{Peter Hollows}
\date{\today}

\begin{document}
\maketitle

\begin{abstract}
Bitcoin mining searches for a block header whose double-SHA-256 hash falls
below a target, a weak-preimage problem solved in practice by brute force.
We ask whether any classical algorithm beats brute force, and answer with a
complete accounting of where an advantage could live. Decomposing SHA-256d
into two layers with opposite properties (a shareable message schedule and
a pseudorandom round function) splits the cost into a per-candidate factor
and a candidate-count factor. We show the per-candidate factor is confined
to three known, bounded regions (midstate caching, the ASICBoost
message-expansion region, and a provably unshareable second hash), so it is
near-fixed. We prove, in the random-oracle model, that the candidate count
cannot fall below $2^{D}$, a bound that survives preprocessing (Bitcoin is
salted) and quantum search (Grover gives only a quadratic speedup), and we
reduce the standard-model claim to a SHA-256 distinguisher. We close the
remaining economic loophole (selecting which header prefixes to mine) by
showing stems are exchangeable under the natural generator, as the round
function's mixing predicts. The result localizes the entire
residual risk to one object: a global algebraic shortcut through SHA-256,
which is a distinguisher. \emph{Mining cannot beat brute force unless
SHA-256 is broken.} Whether that residual is empty is an empirical
question about SHA-256's round function, which this paper isolates but does
not resolve; the public cryptanalytic record is the present evidence.
\end{abstract}

% =====================================================================
\section{The mining search problem}
% =====================================================================

Bitcoin mining is a search problem over an iterated hash. A block header
$H$ is $80$ bytes; the mining hash is the double application
\[
  \Mhash(H) \;=\; \SHA\!\big(\SHA(H)\big),
\]
read as a $256$-bit integer. Mining succeeds when a miner finds an $H$
(within its allowed degrees of freedom: the $32$-bit nonce, the
coinbase extranonce that varies the Merkle root, version bits, and the
timestamp) such that $\Mhash(H) < T$, where $T = 2^{256}\cdot 2^{-D}$
and $D$ is the difficulty expressed in leading zero bits. Equivalently,
$H$ must land in a target set $S \subseteq \{0,1\}^{256}$ of density
$|S|/2^{256} = 2^{-D}$.

The question is operational: \emph{is there any classical algorithm that
finds such an $H$ at materially less expected cost than brute force?}
Brute force tries candidates until one lands in $S$, at expected cost
$\approx 2^{D}$ evaluations. An attack ``beats'' it only if it reduces the
total work by more than a constant factor.

\paragraph{Cost model.}
The atomic unit is one $\SHA$ compression-function call
$h:\{0,1\}^{256}\times\{0,1\}^{512}\to\{0,1\}^{256}$. It applies the
block cipher $E$ (SHACAL-2, keyed by the message $m$) to the chaining
value and feeds the input forward by \emph{wordwise modular addition},
$h(C,m) = C \boxplus E_m(C)$, where $\boxplus$ adds each of the eight
$32$-bit words modulo $2^{32}$. The feed-forward is thus itself a carry
operation, which matters for the round function's resistance to
prediction. One header
evaluation decomposes into exactly three compression calls,
\[
\underbrace{C_1 = h(\mathrm{IV},B_0)}_{\text{block 0 (nonce-independent)}},\quad
\underbrace{C_2 = h(C_1,B_1\Vert\mathrm{pad})}_{\text{block 1 (nonce in word }W_3)},\quad
\underbrace{\mathrm{out}=h(\mathrm{IV},C_2\Vert\mathrm{pad})}_{\text{second hash}},
\]
with $B_0$ the first $64$ header bytes and $B_1$ the remaining $16$ plus
padding. Total mining cost is the number of candidate evaluations
multiplied by the compression calls per candidate that sharing cannot
remove. These two factors are governed by different structure, which the
next section makes precise.

% =====================================================================
\section{Two layers}\label{sec:twolayers}
% =====================================================================

The analysis rests on one structural fact: $\SHAd$ is built from two
kinds of structure with opposite properties.

\begin{center}
\begin{tabular}{ll p{0.36\textwidth}}
\hline
\textbf{layer} & \textbf{property} & \textbf{consequence}\\
\hline
message schedule & fixed recurrence (same per block) & computation is \emph{shareable} across candidates\\
round / compression & pseudorandom & output is \emph{unpredictable}; search is \emph{irreducible}\\
\hline
\end{tabular}
\end{center}

The only avenues to beat brute force live in one layer or the other:
either reduce the \emph{per-candidate cost} by sharing schedule
computation (Section~\ref{sec:amort}), or reduce the \emph{candidate
count} by predicting or skipping evaluations (Section~\ref{sec:search}).
We show the first avenue is bounded to known constant-factor techniques,
and the second is impossible in the ideal model and equivalent to breaking
$\SHA$ in the standard model. Because the two avenues act on different
layers, a gain in one neither requires nor enables a gain in the other,
and the rest of the paper addresses them in turn. \emph{Why} the round
function is pseudorandom (the mechanism behind the second row) is beyond
the scope of this paper; here we take it as the hypothesis to be localized
and reduced.

% =====================================================================
\section{The amortization envelope}\label{sec:amort}
% =====================================================================

Written as a product, the cost is
\[
  \text{cost} \;\ge\;
  \underbrace{(\text{candidate evaluations})}_{\text{locked at }2^D,\ \S\ref{sec:search}}
  \;\times\;
  \underbrace{(\text{compressions per candidate after maximal sharing})}_{\text{this section}}.
\]
The left factor is locked at $2^{D}$ by the search lower bound. Hence any
economically meaningful win against brute force must come from the right
factor: per-candidate computation sharing. We bound it by classifying
which of the three compressions can be shared across \emph{distinct}
candidates.

\begin{itemize}
\item \textbf{Block~0 (midstate).} Nonce-independent, and shared across
  the entire nonce sweep of a fixed header prefix, so its amortized cost
  per candidate falls to zero. This is midstate caching, standard practice
  in deployed miners.
\item \textbf{Second hash.} Effectively unshareable. Two candidates share
  the second-hash compression only if they share a first-hash digest,
  which is a $\SHA$ collision. Finding one costs about $2^{128}$, far more
  than mining the block itself (cost $2^{D}$ with $D < 128$), so the
  second hash cannot be shared profitably and every candidate pays one
  full second-hash compression. This rests on the collision resistance of
  $\SHA$, the mildest cryptographic assumption in the paper, and no
  input-side technique avoids it.
\item \textbf{Block~1 expansion.} Block~1 of the first hash contains the
  nonce, and only its message-schedule expansion is shareable, across
  small batches of candidates engineered to collide in the expansion.
  ASICBoost~\cite{asicboost} exploits exactly this region, covertly via
  the Merkle root or overtly via version rolling (standardized in the
  reserved \texttt{nVersion} bits~\cite{bip320}). It is a deployed
  technique worth about $20\%$, and the only region where a sharing win
  can live.
\end{itemize}

These three regions are the three compressions of the cost model, so they
exhaust the per-candidate work: the sharing envelope is exactly midstate
(at floor), the ASICBoost block-1 expansion (near-saturated), and the
second hash (provably unshareable), with nothing left over. Sharing
\emph{between} the compressions is excluded by the same data dependency:
each compression consumes the previous one's output, so two candidates
share a later compression only if they collide in an earlier digest, which
is again a $\SHA$ collision at cost $2^{128}$. The
amortization avenue is therefore known, bounded, and already occupied,
and the remainder of this paper concerns the candidate-count factor,
where the open question lay.

% =====================================================================
\section{The search lower bound}\label{sec:search}
% =====================================================================

We first establish, unconditionally in the ideal model, that the
candidate-count factor cannot be reduced below $2^{D}$ in expectation;
then we transfer the statement to the standard model by a tight
reduction.

\subsection{Ideal model}

Model the search map $H\mapsto\Mhash(H)$ as a random oracle: with the
compression function idealized, fresh inputs receive uniform, independent
outputs (composition of ideal primitives is not automatic in general; the
construction-level caveats are handled in the standard-model reduction
below). This is the standard idealization under which proof-of-work
hardness is proved \cite{backbone}.

\begin{theorem}[ROM search optimality]\label{thm:rom}
Let $\mathcal{A}$ make at most $q$ queries to the oracle and output a
candidate $H^\ast$. Then
\[
  \Pr[\Mhash(H^\ast)\in S] \;\le\; (q+1)\,2^{-D},
\]
over the oracle and $\mathcal{A}$'s coins. Consequently $\mathcal{A}$
needs $q \ge 2^{D-1}-1$ queries to succeed with probability $\ge \tfrac12$,
and the expected number of queries to the first success is
$\ge 2^{D}$.
\end{theorem}

\begin{proof}[Proof sketch]
Without loss of generality the queries $x_1,\dots,x_q$ are distinct.
Each $\mathrm{RO}(x_i)$ is uniform on $\{0,1\}^{256}$, independent of all
prior responses, so $\Pr[\mathrm{RO}(x_i)\in S]=2^{-D}$. If $H^\ast$ was
queried, success is among these events; if not, $\mathrm{RO}(H^\ast)$ is
uniform and independent, contributing a further $2^{-D}$. A union bound
over the $q$ queries and the final output gives $(q+1)2^{-D}$. The query
statement is immediate. For the expectation, random-oracle freshness makes
each query succeed with probability at most $2^{-D}$ conditioned on all
prior responses, so the first-success time $Q$ stochastically dominates a
$\mathrm{Geometric}(2^{-D})$ variable: $\Pr[Q>q]\ge(1-2^{-D})^{q}$. Hence
$\mathbb{E}[Q]=\sum_{q\ge0}\Pr[Q>q]\ge\sum_{q\ge0}(1-2^{-D})^{q}=2^{D}$.
\end{proof}

Brute force achieves expected $2^{D}$ evaluations, so it is
exactly query-optimal in expectation: no candidate ordering, feature,
or filter reduces the count below $2^{D}$ in this model.

\subsection{Closing the realistic loopholes}

\paragraph{Precomputation and time--memory trade-offs.}
A miner might precompute an enormous table offline. The auxiliary-input
(non-uniform) random-oracle framework~\cite{unruh,cdgs,cdg} models an
adversary with $S$ bits of arbitrary advice and $q$ online queries and
shows that \emph{salting} generically defeats preprocessing. For the
salted inversion problem the advantage is bounded by
$ST/(K\alpha)+T/\alpha$, which for a salt as wide as the function's
output ($K=N$) collapses to $\bigO(T/N)$~\cite{dgk}, i.e.\ the
non-preprocessing bound up to constants. Bitcoin is salted by
construction: each header commits to \texttt{prev\_block\_hash}, a fresh,
unpredictable $256$-bit value, and to the Merkle root of fresh
transactions, so the per-block salt is effectively as wide as the
function output.
Classical time--memory trade-offs~\cite{hellman} and their rainbow-table
refinement~\cite{rainbow} require a \emph{fixed} target function so that
one precomputed table amortizes over all instances; a fresh per-block
salt makes each block a distinct function, and offline computation cannot
target a block whose predecessor is not yet known, so these trade-offs do
not apply.

\paragraph{Quantum.}
Grover search gives a quadratic speedup, $\bigO(2^{D/2})$ queries, and
this is optimal (the BBBV lower bound~\cite{bbbv}): no exponential
speedup exists even quantumly. This carries over to the mining setting
specifically: the post-quantum analysis of the Bitcoin backbone proves
that each quantum query is worth only $\bigO(p^{-1/2})$ classical queries,
with a unified hybrid quantum/classical search bound interpolating the two
regimes~\cite{pqbackbone,hybrid}. The quadratic factor is economically
irrelevant at current scale, since a coherent $\SHAd$ Grover oracle costs
orders of magnitude more per query than a classical ASIC gate-pass.

\paragraph{Standard-model reduction.}
A concrete miner $\mathcal{M}$ against the true $\SHA$ with expected cost
$2^{D}/\alpha$ for super-constant $\alpha$ succeeds at a rate exceeding the
ceiling that Theorem~\ref{thm:rom} imposes on every algorithm in the ROM.
Running $\mathcal{M}$ against a real or an ideal primitive and checking
whether it beats the bound is then a distinguishing experiment, formalized
by indifferentiability~\cite{mrh}: if the mining map is indifferentiable
from a random oracle, the ROM bound transfers to the standard model up to
the distinguishing advantage, and standard-model proof-of-work primitives
of exactly this lower-bound shape have been formalized~\cite{gkp}. One
subtlety must be handled. Plain (strengthened) Merkle--Damg\aa rd is
\emph{not} indifferentiable from a random oracle, because of length
extension: from $\SHA(x)$ one can predict $\SHA(x\Vert\mathrm{pad}\Vert
y)$ without knowing $x$~\cite{coron}. This is the structure that the
\emph{outer} application in $\SHAd=\SHA(\SHA(\cdot))$ is designed to remove
(length extension specifically): re-hashing the full digest is intended to
seal the inner chaining value behind a second, length-committed
invocation, so that the length-extension distinguisher does not reach the
mining map. The literature comes close to this composition from both
sides without covering it. Coron et al.~\cite{coron} prove plain
Merkle--Damg\aa rd indifferentiable (with birthday bounds) on prefix-free
inputs, of which fixed-length inputs are a special case (so a
\emph{single} $\SHA$ restricted to 80-byte headers is covered), and
separately prove indifferentiability of the double hash
$H(H(0^{\kappa}\Vert M))$, which differs from $\SHAd$ only by the
prepended zero block whose role is to domain-separate the outer
single-compression call from adversarially controlled first message
blocks; $\SHAd$ omits that block, so the theorem does not literally
apply. Dodis et al.~\cite{drst} analyze the unrestricted double hash
$H^{2}(M)=H(H(M))$ and show it is indifferentiable only with weak
bounds (any successful simulator must make
$\Omega(\min\{q_{1}q_{2},2^{n/2}\})$ queries, capping
composition-derived security near $2^{n/4}$); their differentiating
attack, however, requires feeding $n$-bit digests back into the
construction, which the mining map's fixed 80-byte domain forbids, and
they prove no positive result for the domain-restricted composition.
No published theorem covers $\SHAd$ on its fixed-length domain, so we
state indifferentiability of the mining map as an explicit assumption,
consistent with formal treatments of Bitcoin that model the mining
function as a random oracle directly~\cite{backbone}.
Indifferentiability composes for single-stage games;
multi-stage settings can require more care~\cite{rss}, but the mining
experiment is single-stage: one adversary, one search.

\begin{theorem}[conditional, informal]\label{thm:reduction}
If the mining map $\Mhash=\SHA(\SHA(\cdot))$ is indifferentiable from a
random oracle (with its compression function modeled as ideal), then no
mining algorithm achieves expected cost below $2^{D}/\bigO(1)$ candidate
evaluations, up to the indifferentiability advantage; an algorithm that
does yields a distinguisher against $\SHA$.
\end{theorem}

\paragraph{The provability wall.}
A fully unconditional, standard-model version of
Theorem~\ref{thm:reduction} is beyond current techniques. Mining for the
fixed function $\SHA$ is finite, so the statement is not literally about
$\mathsf{P}$ versus $\mathsf{NP}$; an unconditional proof would be an
explicit circuit lower bound, which no method can establish for a
function of this complexity. In the asymptotic framing, for a scaled
family of $\SHA$-like functions, proving that search requires
$\sim 2^{D}$ work would exhibit an explicit one-way function and so
separate $\mathsf{P}$ from $\mathsf{NP}$. Either way, the conditional
reduction above is the strongest conclusive statement available, and it
localizes the entire risk to one assumption.

\subsection{Localization}\label{sec:localization}

\begin{corollary}\label{cor:localization}
Every economically meaningful win against brute force is confined to the
per-candidate sharing factor of Section~\ref{sec:amort}. Output
prediction and search-skipping are impossible in the ideal model
(Theorem~\ref{thm:rom}) and equivalent to a $\SHA$ break in the standard
model (Theorem~\ref{thm:reduction}).
\end{corollary}

Whether the round-function layer admits such prediction is exactly the
residual this localization isolates: an empirical question about $\SHA$,
not one the reduction settles. One economic loophole in the corollary
remains to be closed directly, and we do so next.

% =====================================================================
\section{Stem selection}\label{sec:stem}
% =====================================================================

Theorem~\ref{thm:rom} locks the candidate count for a \emph{fixed} header
prefix. A miner with freedom over the prefix (the extranonce, version
bits, and timestamp) might still hope that some prefixes, which we call
\emph{stems}, yield more readily than others, and allocate effort to the
best. This would not contradict the per-stem bound, but it would beat
brute force in aggregate. We close the loophole: under the natural
generator, stems are exchangeable, and any yield bias is bounded by the
round function's mixing.

Define the yield of a stem as the extreme-value statistic of a full nonce
sweep. Sweeping $K$ nonces and recording the best (smallest) hash, set
$Y=-\log_2(\text{min-hash}/2^{256})$, the number of leading zero bits of
the best hash found. Under the random-oracle idealization the per-nonce
hashes are i.i.d.\ uniform, so $Y$ follows a Gumbel law with
$\mathbb{E}[Y]\approx\log_2 K + 0.83$ and variance $\approx 3.41$,
\emph{independent of the stem}. A stem carries exploitable structure only
if its yield deviates from this law, equivalently if some cheap,
nonce-independent partition of stems has a heavier tail than the null.

We tested this directly. Under a synthetic mining-template generator,
sweeping on the order of $10^{4}$ nonces over each of $\sim\!3\times10^{4}$
stems, no partition drawn from a pre-registered feature catalogue beat its
null tail-hit rate through a three-stage discovery / confirmation /
fresh-holdout protocol with multiple-testing correction (no partition of
several dozen survived across hundreds of buckets); a positive control
with an injected per-partition bias was recovered, and a separate family of
synthetic-uniform negative controls produced no lift. Under the tested generator and
catalogue, stems carry no exploitable yield. (Full pre-registration,
generator, feature catalogue, and per-bucket results accompany this
paper as code. The closure is for the synthetic generator; we did not
obtain operational pool work-queue data, so the claim is the natural-model
statement, not an audit of any specific deployed generator.)

The measurement is what one expects if the round function thoroughly mixes
the stem bits into the output. A stem fixes some input bits, but their
influence on any one output bit is diffused through the many carry-coupled
modular additions between the header and the digest, so no cheap,
nonce-independent feature of a stem survives to bias a nonce sweep. The
empirical exchangeability is the evidence; this mixing picture only
indicates why it is unsurprising.

% =====================================================================
\section{Related work and conclusion}\label{sec:related}
% =====================================================================

Proof-of-work security is conventionally proved in the random-oracle
model~\cite{backbone}, with standard-model~\cite{gkp} and
post-quantum~\cite{pqbackbone,hybrid} analyses extending it. As we make
explicit, the impossibility of beating brute force is not provable
unconditionally in the standard model without resolving open hardness
questions. The contribution here is the \emph{complete accounting}. A
bounded sharing envelope, an unconditional ideal-model search floor, a
tight standard-model reduction, and a closed stem-selection loophole
together localize the entire residual to one object: a global algebraic
shortcut through $\SHA$ that resolves the target output without resolving
the carries on the path, which by the reduction is a distinguisher against
$\SHA$. ASICBoost~\cite{asicboost} is the canonical
input-side win, and it occupies the only shareable region we identify.
Auxiliary-input bounds~\cite{cdgs,cdg,unruh,dgk} and the BBBV quantum lower
bound~\cite{bbbv} close the preprocessing and quantum loopholes. The public
cryptanalytic record is consistent with this ceiling: reduced-round
collision attacks reach only $31$ of $64$ steps (made practical
in~\cite{lll24}, building on~\cite{mendel13}) and $39$ steps
semi-free-start~\cite{lll24}, and preimages reach $43$ steps at cost
$2^{254.9}$~\cite{aoki09}, all far from the full function.

Whether that residual is empty is an empirical question about $\SHA$'s
round function. The public cryptanalytic record is the present evidence:
no shortcut of the required kind has been found, and the reduced-round
results above remain far from the full function.

SHA-256d mining cannot beat brute force by any route we or the
cryptanalytic record can find, and cannot at all unless $\SHA$ is broken.
The only standing economic wins are the known constant-factor sharing
techniques, and they are bounded.

% =====================================================================
\begin{thebibliography}{99}

\bibitem{backbone}
J.~Garay, A.~Kiayias, N.~Leonardos.
\emph{The Bitcoin Backbone Protocol: Analysis and Applications.}
EUROCRYPT 2015.

\bibitem{asicboost}
T.~Hanke.
\emph{AsicBoost: A Speedup for Bitcoin Mining.}
2016, \texttt{arXiv:1604.00575}.

\bibitem{bip320}
\emph{nVersion Bits for General Purpose Use (BIP~320).}
Bitcoin Improvement Proposals, 2018.

\bibitem{unruh}
D.~Unruh.
\emph{Random Oracles and Auxiliary Input.}
CRYPTO 2007.

\bibitem{cdgs}
S.~Coretti, Y.~Dodis, S.~Guo, J.~Steinberger.
\emph{Random Oracles and Non-Uniformity.}
EUROCRYPT 2018.

\bibitem{cdg}
S.~Coretti, Y.~Dodis, S.~Guo.
\emph{Non-Uniform Bounds in the Random-Permutation, Ideal-Cipher, and
Generic-Group Models.}
CRYPTO 2018.

\bibitem{dgk}
Y.~Dodis, S.~Guo, J.~Katz.
\emph{Fixing Cracks in the Concrete: Random Oracles with Auxiliary Input,
Revisited.}
EUROCRYPT 2017.

\bibitem{hellman}
M.~E.~Hellman.
\emph{A Cryptanalytic Time-Memory Trade-Off.}
IEEE Trans.\ Inf.\ Theory 26(4), 1980.

\bibitem{rainbow}
P.~Oechslin.
\emph{Making a Faster Cryptanalytic Time-Memory Trade-Off.}
CRYPTO 2003.

\bibitem{bbbv}
C.~Bennett, E.~Bernstein, G.~Brassard, U.~Vazirani.
\emph{Strengths and Weaknesses of Quantum Computing.}
SIAM J.\ Comput.\ 26(5), 1997.

\bibitem{pqbackbone}
A.~Cojocaru, J.~Garay, A.~Kiayias, F.~Song, P.~Wallden.
\emph{Quantum Multi-Solution Bernoulli Search with Applications to
Bitcoin's Post-Quantum Security.}
Quantum 7:944, 2023.

\bibitem{hybrid}
A.~Cojocaru, J.~Garay, F.~Song.
\emph{Generalized Hybrid Search with Applications to Blockchains and Hash
Function Security.}
ASIACRYPT 2024.

\bibitem{mrh}
U.~Maurer, R.~Renner, C.~Holenstein.
\emph{Indifferentiability, Impossibility Results on Reductions, and
Applications to the Random Oracle Methodology.}
TCC 2004.

\bibitem{gkp}
J.~Garay, A.~Kiayias, G.~Panagiotakos.
\emph{Proofs of Work for Blockchain Protocols.}
IACR Cryptology ePrint Archive, Report 2017/775, 2017.

\bibitem{coron}
J.-S.~Coron, Y.~Dodis, C.~Malinaud, P.~Puniya.
\emph{Merkle-Damg\aa rd Revisited: How to Construct a Hash Function.}
CRYPTO 2005.

\bibitem{drst}
Y.~Dodis, T.~Ristenpart, J.~Steinberger, S.~Tessaro.
\emph{To Hash or Not to Hash Again? (In)differentiability Results for
$H^2$ and HMAC.}
CRYPTO 2012.

\bibitem{rss}
T.~Ristenpart, H.~Shacham, T.~Shrimpton.
\emph{Careful with Composition: Limitations of the Indifferentiability
Framework.}
EUROCRYPT 2011.

\bibitem{lll24}
Y.~Li, F.~Liu, G.~Wang.
\emph{New Records in Collision Attacks on SHA-2.}
EUROCRYPT 2024.

\bibitem{mendel13}
F.~Mendel, T.~Nad, M.~Schl\"affer.
\emph{Improving Local Collisions: New Attacks on Reduced SHA-256.}
EUROCRYPT 2013.

\bibitem{aoki09}
K.~Aoki, J.~Guo, K.~Matusiewicz, Y.~Sasaki, L.~Wang.
\emph{Preimages for Step-Reduced SHA-2.}
ASIACRYPT 2009.

\end{thebibliography}

\end{document}
